\section{Security}

\subsection{Security Assessment}

\subsubsection{Risk Identification}

\paragraph{Assets}
\begin{itemize}[noitemsep]
    \item \textbf{Web application}
    \item \textbf{User data} in PostgreSQL
    \item \textbf{Server}
    \item \textbf{Secrets}
    \item \textbf{CI/CD pipeline}
    \item \textbf{Monitoring stack}
\end{itemize}

\paragraph{Threat sources}
\begin{itemize}[noitemsep]
    \item External attackers
    \item Malicious users
    \item Supply-chain compromise (Go modules, base images, GitHub Actions)
    \item Developer mistakes
    \item Cloud / hardware failure
\end{itemize}

\paragraph{Risk scenarios}

\begin{itemize}[noitemsep]
    \item \textbf{R1}: Attacker brute-forces \texttt{POST /login} to take over user accounts.
    \item \textbf{R2}: Attacker uses the hardcoded simulator API token in \texttt{internal/web/router.go} to post or read messages as any user.
    \item \textbf{R3}: Attacker scrapes or floods the API (\texttt{/msgs?no=\ldots}) and takes down the server.
    \item \textbf{R4}: Hetzner server is lost or compromised.
    \item \textbf{R5}: A third-party GitHub Action used in CI is compromised and obtains \texttt{SSH\_PRIVATE\_KEY} / \texttt{SERVER\_IP}.
    \item \textbf{R6}: Attacker can abuse the exposed internal services host ports in \texttt{docker-compose.yml}.
    \item \textbf{R7}: Attacker brute-forces our short passwords.
\end{itemize}

\subsubsection{Risk Analysis}

\paragraph{Likelihood \& Impact}

\subparagraph{Scales}
\begin{itemize}[noitemsep]
    \item \textbf{Likelihood:} 1 (Rare), 2, 3, 4, 5 (Almost certain)
    \item \textbf{Impact:} 1 (Negligible), 2, 3, 4, 5 (Severe)
\end{itemize}

\begin{center}
\begin{tabular}{lccc}
\toprule
\textbf{Risk scenario} & \textbf{Likelihood} & \textbf{Impact} & \textbf{Score} \\
\midrule
R1  & 4 & 4 & \textbf{16} \\
R2  & 5 & 3 & \textbf{15} \\
R6  & 5 & 3 & \textbf{15} \\
R4  & 3 & 5 & \textbf{15} \\
R7 & 3 & 4 & \textbf{12} \\
R5  & 2 & 5 & \textbf{10} \\
R3  & 3 & 3 & \textbf{9} \\
\bottomrule
\end{tabular}
\end{center}

\paragraph{Risk Matrix}

\begin{center}
\renewcommand{\arraystretch}{1.3}
\begin{tabular}{c|ccccc}
\multicolumn{1}{c}{} & \multicolumn{5}{c}{\textbf{Impact}} \\
\textbf{Likelihood} & 1 & 2 & 3 & 4 & 5 \\
\hline
5 & --- & --- & R2, R6   & ---  & ---     \\
4 & --- & --- & ---      & R1   & ---     \\
3 & --- & --- & R3       & R7  & R4      \\
2 & --- & --- & ---      & ---  & R5      \\
1 & --- & --- & ---      & ---  & ---     \\
\end{tabular}
\end{center}

\paragraph{Addressing the issues}

\subparagraph{R2}
We have moved the hardcoded simulator API token out of \texttt{internal/web/router.go} and into \texttt{.env}.

\subparagraph{R6}
We removed the ports exposures for Grafana and Kibana from texttt{docker-compose.yml} -- the services remain reachable only through Caddy on \texttt{grafana.kulturbase.dk} and \texttt{kibana.kulturbase.dk}.

\subparagraph{R7}
For R10, we replaced every service credential in \texttt{.env} with new updated credentials random strings generated using \texttt{openssl}.

We prioritized R2 and R6 because while they were critical, they were also relatively straightforward to fix. R7 was tackled since it was also a simple fix that was long overdue.

\subsection{Firewall}

We use both Hetzner's built-in firewall and \texttt{ufw} on the server. Only ports \texttt{80}, \texttt{443}, and \texttt{22} are open to incoming traffic.

\subsection{TLS}

We needed HTTPS in front of the app, and the lecture recommended using Nginx with Certbot. We chose Caddy instead because it combines both into a single tool: it automatically requests and renews Let's Encrypt certificates and redirects HTTP traffic to HTTPS. Our TLS setup is a single Caddyfile that routes \texttt{kulturbase.dk} and the \texttt{grafana.} and \texttt{kibana.} subdomains to their respective containers.

\subsubsection{Container Hardening}

\subsubsection{Not running as root}

We have added a \texttt{USER} directive to our Dockerfile to prevent the images running as \texttt{root}. We have additionally run \texttt{docker scout} on our images which revealed no critical vulnerabilities. Our base images are up-to-date with versions \texttt{golang:1.25} and \texttt{alpine:3.23}, neither of which has any known vulnerabilities.

\subsubsection{Image scan for vulnerabilities}

The initial scan of our \texttt{minitwit} image using \texttt{docker scout cves --only-severity critical,high} revealed 4 vulnerabilities: 2 critical in the \texttt{github.com/jackc/pgx/v5} Go dependency and 2 high in the Alpine 3.19 image. The following fixes were applied:
\begin{itemize}[noitemsep]
    \item Updated \texttt{pgx/v5} from \texttt{v5.6.0} to \texttt{v5.9.2}
    \item Updated \texttt{alpine:3.19} to \texttt{alpine:3.23}
\end{itemize}

A follow-up scan reports 0 critical and 0 high risk vulnerabilities.

\subsection{Enhance your CI Pipelines with multiple security related static analysis tools}

CI runs CodeQL on every push and PR to \texttt{main}. The deploy job depends on \texttt{codeql}, and branch protection requires the code scanning check to pass before merge, so any finding above a specific threshold blocks merge and deployment. 

CI also runs Trivy on every push and PR to \texttt{main}: the \texttt{image\_scan} job builds the container and scans it for fixable HIGH/CRITICAL vulnerabilities in OS packages and Go modules. The deploy job depends on \texttt{image\_scan}, so any vulnerable image is blocked from deployment.
